\section{Related work}\label{sec:related}
\paragraph{Raw-event attacks and constrained retiming.}
Surrogate-gradient learning provides a practical route through the nondifferentiable dynamics of spiking networks \cite{neftci2019}. B\"uchel et al. study event additions and removals in spiking vision models \cite{buchel2022attacks}, a different feasible set from identity-preserving retiming. An attack on an aggregated tensor and an attack on the underlying event stream need not have the same feasible set. Yao et al. attack raw event data using a probabilistic treatment of discrete event values and sparsity control \cite{yao2024rawevent}. Their construction permits event-value and support changes, whereas our native-event threat model preserves event identities, coordinates, polarities and count. Du et al. investigate raw-event attacks with configurable latencies \cite{du2026latencies}, providing a particularly relevant connection to temporal manipulation. These works motivate attacks at the event interface without making their constraints interchangeable with ours.

\paragraph{Capacity constraints and representation-specific budgets.}
Yu et al. formalize capacity-one spike retiming and budgets on maximum displacement, total delay and the number of retimed spikes \cite{yu2026}. Our classical packet-grid comparison uses a local implementation of this framework, with completed results for maximum displacement and retimed-packet count. The native-timestamp experiment instead permits multiple events to contribute to a frame cell. Its duration-scaled timestamp radius cannot be identified with a one-bin packet displacement. We consequently compare attacks within each stated regime and do not rank the regimes by their raw ASR values.

\paragraph{Hybrid spiking and quantum models.}
Quantum temporal models encompass distinct computational constructions. Brand and Petruccione develop a quantum leaky integrate-and-fire neuron and network \cite{brand2024}; Jha et al. combine SNN-derived EEG features with quantum-kernel classification \cite{jha2025hybrid}; and SPATE studies spike-derived temporal quantum encodings \cite{spate2026}. QSNN-v3 instead composes classical Conv/LIF dynamics, a learned angle map, a data-re-uploading variational circuit \cite{perez2020}, and a linear classifier on the complete measurement-probability vector. These differences concern the model and interface, so accuracy or security claims do not transfer directly between the constructions.

\paragraph{Adversarial robustness of quantum classifiers.}
Wendlinger et al. compare quantum and classical classifiers using transfer attacks, perturbation structure and Lipschitz-based analysis \cite{wendlinger2024comparative}. Gong et al. investigate randomized encodings under specified adversarial-access assumptions \cite{gong2024randomized}, while Dowling et al. establish robustness guarantees for particular quantum-classifier settings \cite{dowling2026guarantees}. Khatun and Usman use classical autoencoders to distill quantum adversarial manipulations on image benchmarks \cite{khatun2025distillation}; this is neither an implemented defense here nor evidence for temporal-retiming protection. These studies motivate stating the encoding and adversary explicitly. Their guarantees and defenses cannot be assigned to QSNN-v3 without checking the corresponding assumptions, especially because our perturbation enters before a classical event-processing frontend.

\paragraph{Evaluation strength and computational controls.}
PGD is a standard iterative baseline \cite{madry2018}, but the failure of one optimizer is insufficient evidence of robustness \cite{carlini2019,athalye2018obfuscated}. Diverse attacks and improved optimization can expose failures missed by an individual configuration \cite{croce2020reliable}. We therefore report paired outcomes under candidate, forward-pass and elapsed-time controls. These controls answer different resource questions and do not establish optimality of either attack. Standard image-space attack ensembles also require adaptation before they respect our order, timing and packet constraints.


The supplementary conceptual-positioning table distinguishes these interfaces without ranking external ASRs.
